
\documentclass{article}
\usepackage{colortbl}
\usepackage{makecell}
\usepackage{multirow}
\usepackage{supertabular}

\begin{document}

\newcounter{utterance}

\centering \large Interaction Transcript for game `dond', experiment `coop\_en', episode 12 with reflexion{-}Qwen3.5{-}9B.
\vspace{24pt}

{ \footnotesize  \setcounter{utterance}{1}
\setlength{\tabcolsep}{0pt}
\begin{supertabular}{c@{$\;$}|p{.15\linewidth}@{}p{.15\linewidth}p{.15\linewidth}p{.15\linewidth}p{.15\linewidth}p{.15\linewidth}}
    \# & $\;$A & \multicolumn{4}{c}{Game Master} & $\;\:$B\\
    \hline

    \theutterance \stepcounter{utterance}  
    & & \multicolumn{4}{p{0.6\linewidth}}{
        \cellcolor[rgb]{0.9,0.9,0.9}{
            \makecell[{{p{\linewidth}}}]{
                \texttt{\tiny{[P1$\langle$GM]}}
                \texttt{You have attempted this same task before and did not succeed. These are the plans you wrote after those attempts.} \\
\texttt{Trial \#1: New plan: I failed because I included explanatory text and conversational phrases ("Message 1:", "What do you think?", "Let's try to be more direct") alongside my proposal, which the Game Master interpreted as a secret proposal due to the presence of square brackets or simply rejected as non{-}compliant with the "no response to rules" instruction. The rules explicitly state that messages are directly transmitted and that using square brackets triggers a secret proposal. To succeed, I must send a single, clean message containing only the negotiation text without any labels, brackets, or meta{-}commentary. I will draft a message that clearly states my desired split (1 envelope, 1 spoon for me; 1 spoon, 2 shells, 1 pepper for you) in plain text, ensuring no square brackets are used unless I am certain I am making a final secret proposal (which I should avoid until the negotiation phase is exhausted). I will strictly omit any introductory phrases like "Message 1:" or "I suggest".} \\
\\ 
\texttt{Trial \#2: New plan: I failed because I submitted a secret proposal (`{[}Proposal: 1 pepper, 2 shells{]}`) that was incomplete and did not specify the allocation for the other player (the Game Master). The rules state that proposals must be complementary to succeed; by only listing items for myself, the system could not verify if the remaining items (2 spoons, 1 envelope) were sufficient and correctly assigned to the other player to form a valid, complementary deal. To succeed, I must ensure my secret proposal explicitly defines the full split: clearly stating which items I want and which items the other player gets, ensuring the sum of requested items equals the total available inventory (2 shells, 2 spoons, 1 envelope, 1 pepper). I will draft a proposal like `{[}Proposal: 1 pepper, 2 shells for me; 2 spoons, 1 envelope for you{]}` to guarantee complementarity.} \\
\\ 
\texttt{Trial \#3: New plan: I failed because I violated the critical rule that square brackets `{[}{]}` trigger an immediate secret proposal, which the Game Master rejected as an invalid move during the open negotiation phase. The rules explicitly state that using square brackets in a message to the other player will be interpreted as a secret proposal, and since I was not supposed to make a secret proposal yet (I had 5 messages available), this caused the abort. To succeed, I must communicate my desired split using only plain text without any square brackets, ensuring the allocation is clear and complementary (e.g., "I want the envelope and one spoon; you take the pepper, the second spoon, and both shells"). I will strictly avoid any formatting that resembles the secret proposal syntax until I have exhausted my open negotiation messages or intentionally choose to end the game.} \\
\\ 
\\ 
\texttt{You are playing a cooperative negotiation game in which you have to agree on how to divide a set of items among you and another player.} \\
\\ 
\texttt{Rules:} \\
\texttt{(a) You and the other player are given a set of items. Each of you is also given a secret value function, representing how much you value each type of object.} \\
\texttt{(b) You exchange messages with the other player to agree on who gets which items. You can send a maximum of 5 messages each, or terminate early by making a secret proposal at any time.} \\
\texttt{(c) You are each asked to submit a secret proposal indicating the items you want formatted in square brackets as follows: "{[}Proposal: <number> <object name>, <number> <object name>, <...>{]}"} \\
\texttt{(d) If your proposals are complementary, i.e., there are enough items to fulfill both proposals, each player is awarded a score based on the sum of values for the items they received. Otherwise, both of you get zero points.} \\
\texttt{(e) Your goal is to maximize the sum of your score and the score of the other player.} \\
\\ 
\texttt{Let us start.} \\
\\ 
\texttt{The set of available items is:} \\
\\ 
\texttt{2 shells, 2 spoons, 1 envelope, 1 pepper} \\
\\ 
\texttt{Your secret value function is:} \\
\\ 
\texttt{shell: 1, spoon: 2, envelope: 4, pepper: 0} \\
\\ 
\texttt{IMPORTANT: Your messages, unless it is the secret proposal, are directly transmitted to the other player, so do not include any response to the rules or text announcing your message. To make a secret proposal, use the indicated format. Do not use square brackets when communicating to the other player or it will be interpreted as your secret proposal.} \\
            }
        }
    }
    & \\ \\

    \theutterance \stepcounter{utterance}  
    & \multicolumn{4}{p{0.6\linewidth}}{
        \cellcolor[rgb]{0.9,0.9,0.9}{
            \makecell[{{p{\linewidth}}}]{
                \texttt{\tiny{[P1$\rangle$GM]}}
                \texttt{I value the envelope most, so I propose we split the items this way: I take the envelope and one spoon, and you take the pepper, the second spoon, and both shells. This gives us the envelope (worth 4 to me) and the shells (worth 2 to you), maximizing our combined score.} \\
            }
        }
    }
    & & \\ \\

    \theutterance \stepcounter{utterance}  
    & & & \multicolumn{2}{p{0.3\linewidth}}{
        \cellcolor[rgb]{0.9,0.9,0.9}{
            \makecell[{{p{\linewidth}}}]{
                \texttt{\tiny{[GM$|$GM]}}
                \texttt{continue} \\
            }
        }
    }
    & & \\ \\

    \theutterance \stepcounter{utterance}  
    & & & \multicolumn{2}{p{0.3\linewidth}}{
        \cellcolor[rgb]{0.9,0.9,0.9}{
            \makecell[{{p{\linewidth}}}]{
                \texttt{\tiny{[GM$|$GM]}}
                \texttt{I value the envelope most, so I propose we split the items this way: I take the envelope and one spoon, and you take the pepper, the second spoon, and both shells. This gives us the envelope (worth 4 to me) and the shells (worth 2 to you), maximizing our combined score.} \\
            }
        }
    }
    & & \\ \\

    \theutterance \stepcounter{utterance}  
    & & \multicolumn{4}{p{0.6\linewidth}}{
        \cellcolor[rgb]{0.9,0.9,0.9}{
            \makecell[{{p{\linewidth}}}]{
                \texttt{\tiny{[GM$\rangle$P2]}}
                \texttt{You have attempted this same task before and did not succeed. These are the plans you wrote after those attempts.} \\
\texttt{Trial \#1: New plan: I failed because I spent my turns analyzing the math and explaining the error instead of immediately submitting the required secret proposal format. The game rules state that to finalize the deal, I must use the exact string `{[}Proposal: <items>{]}`. In my last turn, I wrote the proposal but did not include the opening square bracket `{[}`, causing the system to reject it as a valid secret proposal. Next time, upon identifying the optimal allocation (Me: 1 pepper, 2 shells; You: 2 spoons, 1 envelope), I will immediately output the exact string `{[}Proposal: 1 pepper, 2 shells{]}` without any preceding text or analysis to ensure the game terminates successfully with the maximum combined score.} \\
\\ 
\texttt{Trial \#2: New plan: I failed because I included explanatory text before the secret proposal, which the rules explicitly forbid ("Do not include any response to the rules or text announcing your message"). The system interpreted the entire message as a secret proposal, but the format was invalid due to the extra text. Next time, upon agreeing on the split (Me: 1 pepper, 2 shells; You: 2 spoons, 1 envelope), I will output **only** the exact string `{[}Proposal: 1 pepper, 2 shells{]}` with no preceding explanation, analysis, or conversational filler to ensure the game terminates successfully.} \\
\\ 
\\ 
\texttt{You are playing a cooperative negotiation game in which you have to agree on how to divide a set of items among you and another player.} \\
\\ 
\texttt{Rules:} \\
\texttt{(a) You and the other player are given a set of items. Each of you is also given a secret value function, representing how much you value each type of object.} \\
\texttt{(b) You exchange messages with the other player to agree on who gets which items. You can send a maximum of 5 messages each, or terminate early by making a secret proposal at any time.} \\
\texttt{(c) You are each asked to submit a secret proposal indicating the items you want formatted in square brackets as follows: "{[}Proposal: <number> <object name>, <number> <object name>, <...>{]}"} \\
\texttt{(d) If your proposals are complementary, i.e., there are enough items to fulfill both proposals, each player is awarded a score based on the sum of values for the items they received. Otherwise, both of you get zero points.} \\
\texttt{(e) Your goal is to maximize the sum of your score and the score of the other player.} \\
\\ 
\texttt{Let us start.} \\
\\ 
\texttt{The set of available items is:} \\
\\ 
\texttt{2 shells, 2 spoons, 1 envelope, 1 pepper} \\
\\ 
\texttt{Your secret value function is:} \\
\\ 
\texttt{shell: 3, spoon: 0, envelope: 0, pepper: 4} \\
\\ 
\texttt{IMPORTANT: Your messages, unless it is the secret proposal, are directly transmitted to the other player, so do not include any response to the rules or text announcing your message. To make a secret proposal, use the indicated format. Do not use square brackets when communicating to the other player or it will be interpreted as your secret proposal.} \\
\\ 
\\ 
\texttt{I value the envelope most, so I propose we split the items this way: I take the envelope and one spoon, and you take the pepper, the second spoon, and both shells. This gives us the envelope (worth 4 to me) and the shells (worth 2 to you), maximizing our combined score.} \\
            }
        }
    }
    & \\ \\

    \theutterance \stepcounter{utterance}  
    & & & \multicolumn{4}{p{0.6\linewidth}}{
        \cellcolor[rgb]{0.9,0.9,0.9}{
            \makecell[{{p{\linewidth}}}]{
                \texttt{\tiny{[GM$\langle$P2]}}
                \texttt{{[}Proposal: 1 pepper, 2 shells{]}} \\
            }
        }
    }
     \\ \\

    \theutterance \stepcounter{utterance}  
    & & & \multicolumn{2}{p{0.3\linewidth}}{
        \cellcolor[rgb]{0.9,0.9,0.9}{
            \makecell[{{p{\linewidth}}}]{
                \texttt{\tiny{[GM$|$GM]}}
                \texttt{proposal} \\
            }
        }
    }
    & & \\ \\

    \theutterance \stepcounter{utterance}  
    & & & \multicolumn{2}{p{0.3\linewidth}}{
        \cellcolor[rgb]{0.9,0.9,0.9}{
            \makecell[{{p{\linewidth}}}]{
                \texttt{\tiny{[GM$|$GM]}}
                [2, 0, 0, 1]
            }
        }
    }
    & & \\ \\

    \theutterance \stepcounter{utterance}  
    & & \multicolumn{4}{p{0.6\linewidth}}{
        \cellcolor[rgb]{0.9,0.9,0.9}{
            \makecell[{{p{\linewidth}}}]{
                \texttt{\tiny{[P1$\langle$GM]}}
                \texttt{The other player has submitted a secret proposal. It is now your turn to submit a proposal of your own.} \\
            }
        }
    }
    & \\ \\

    \theutterance \stepcounter{utterance}  
    & \multicolumn{4}{p{0.6\linewidth}}{
        \cellcolor[rgb]{0.9,0.9,0.9}{
            \makecell[{{p{\linewidth}}}]{
                \texttt{\tiny{[P1$\rangle$GM]}}
                \texttt{{[}Proposal: 1 envelope, 1 spoon{]}} \\
            }
        }
    }
    & & \\ \\

    \theutterance \stepcounter{utterance}  
    & & & \multicolumn{2}{p{0.3\linewidth}}{
        \cellcolor[rgb]{0.9,0.9,0.9}{
            \makecell[{{p{\linewidth}}}]{
                \texttt{\tiny{[GM$|$GM]}}
                \texttt{proposal} \\
            }
        }
    }
    & & \\ \\

    \theutterance \stepcounter{utterance}  
    & & & \multicolumn{2}{p{0.3\linewidth}}{
        \cellcolor[rgb]{0.9,0.9,0.9}{
            \makecell[{{p{\linewidth}}}]{
                \texttt{\tiny{[GM$|$GM]}}
                [0, 1, 1, 0]
            }
        }
    }
    & & \\ \\

    \theutterance \stepcounter{utterance}  
    & & & \multicolumn{2}{p{0.3\linewidth}}{
        \cellcolor[rgb]{0.9,0.9,0.9}{
            \makecell[{{p{\linewidth}}}]{
                \texttt{\tiny{[GM$|$GM]}}
                [[0, 1, 1, 0], [2, 0, 0, 1]]
            }
        }
    }
    & & \\ \\

\end{supertabular}
}

\end{document}
